\renewcommand{\Anno}{2004}

\renewcommand{\NumeroProblema}{1}

\renewcommand{\Testo}{Determinare tutte le coppie $(m,n)$ di interi
positivi per cui
$$\sqrt[60]{m^{n^{5}-n}}$$
risulta a sua volta un intero.}

\renewcommand{\Risposta}{Sono tutte le coppie in cui $m$ è un
intero positivo qualunque ed $n$ è un intero positivo del tipo $4k$, $4k+1$ o $4k+3$, più tutte le coppie in cui $m=a^{2}$ è
un quadrato perfetto ed $n$ è del tipo $4k+2$.}

\renewcommand{\Aiutino}{Il termine $n^{5}-n=n(n+1)(n-1)(n^{2}+1)$
è sempre pari, multiplo di 3 e di 5, dunque multiplo di 30.  Tale termine è anche multiplo di 4 (dunque di 60) tranne nel caso in
cui $n=4k+2$.}

\renewcommand{\SoluzioneUfficiale}{\TitoloSoluzione Fattorizziamo 
l'esponente
$$n^{5}-n=n(n^{4}-1)=n(n-1)(n+1)(n^{2}+1).$$

Osserviamo ora che tale espressione 
\begin{itemize}
	\item  è sempre divisibile per 2, in quanto o $n$ o $n-1$ sono pari;

	\item  è sempre divisibile per 3, in quanto uno dei 3 fattori consecutivi $n-1$, $n$, $n+1$ è sempre multiplo di 3;

	\item  è sempre divisibile per 5, in quanto se la scrittura decimale di $n$ termina con le cifre 0, 1, 4, 5, 6, 9, allora uno dei primi 3 fattori è multiplo di 5, se invece $n$ termina con 2, 3, 7, 8, allora $n^{2}$ termina con 4 o 9 e quindi $n^{2}+1$ è multiplo di 5;

	\item  è divisibile per 4 se e solo se $n$ non è della forma $4k+2$: infatti se $n$ è del tipo $4k$, $4k+1$ o $4k+3$, allora uno dei 3 fattori consecutivi $n-1$, $n$, $n+1$ è multiplo di 4, se invece $n=4k+2$, allora dei 4 fattori l'unico pari è $n$, che però non è multiplo di 4.
\end{itemize}

In conclusione
\begin{itemize}
	\item  se $n$ è della forma $4k$, $4k+1$ o $4k+3$, allora $n^{5}-n$ è multiplo di 60, quindi l'espressione data rappresenta un intero per ogni scelta dell'intero positivo $m$;

	\item  se $n$ è della forma $4k+2$, allora $n^{5}-n$ è multiplo di 30 ma non di 60, quindi l'espressione data sarà un intero se e solo se $m$ è un quadrato perfetto.
\end{itemize}

Per dimostrare formalmente l'ultimo punto possiamo scrivere $n^{5}-n$ nella forma $30d$, dove $d$ è dispari (altrimenti $n^{5}-n$ sarebbe multiplo anche di 60). Siano ora $m$ ed $h$ interi tali che
$$\sqrt[60]{m^{30d}}=h.$$

Elevando al quadrato otteniamo che
$$m^{d}=h^{2}.$$

Sia ora $p$ un primo che divide $m$. Allora $p$ divide anche $h$, e detti $\alpha$ e $\beta$ gli esponenti con cui $p$ compare nella scomposizione di $m$ ed $h$, rispettivamente, si avrà che
$$\alpha d=2\beta.$$

Pertanto $\alpha d$ è pari e quindi, essendo $d$ dispari, l'unica
possibilità è che $\alpha$ sia pari. Abbiamo cos\`{\i}
dimostrato che nella scomposizione di $m$ tutti i fattori primi compaiono con esponente pari, e da questo segue che $m$ è un quadrato perfetto.

È poi immediato verificare che tale condizione è pure
sufficiente.}

\renewcommand{\SoluzioniAlternative}{
}

\renewcommand{\Osservazioni}{\TitoloOsservazione Abbiamo dimostrato
che $n^{5}-n$ è sempre multiplo di 5 con un ragionamento elementare basato sull'ultima cifra della rappresentazione in base 10 di $n$.  Vale però la pena di ricordare che questo fatto è semplicemente il caso $p=5$ del \emph{piccolo Teorema di Fermat}. Analogamente, tutte le affermazioni che abbiamo fatto su $n^{5}-n$ si possono dimostrare sostituendo ad $n$ i rappresentanti delle classi di
congruenza modulo 3, 4, 5.  }

\renewcommand{\Commenti}{
}